\chapter{First principles: the energy of moving a bit}
\label{ch:firstprinciples}

Before any device, modulation format, or standard, one quantity governs
short-reach interconnect design: the energy required to move a bit from one place
to another. David Miller's work on optical interconnects to silicon lays out the
clearest first-principles framework for this,\sidenote{\citehref{miller2009}{D.~A.~B.~Miller, ``Device
requirements for optical interconnects to silicon chips,'' \emph{Proc. IEEE},
2009}; and \citehref{miller2017}{``Attojoule optoelectronics for low-energy information processing and communications,'' \emph{JLT}, 2017}.} and although the specific numbers are years
old, the scaling arguments are what matter, and they still decide where the optics
go.

\section{The electrical baseline: charging a wire}

To send a bit down an electrical line you charge and discharge its capacitance
through a voltage swing. To order of magnitude,
\[
  E_{\text{elec}} \approx \tfrac{1}{2}\,C\,V^{2},
\]
and the line capacitance grows with length, $C \approx c'\,L$ (with $c'$ on the
order of a hundred-odd femtofarads per millimeter, depending on the medium).%
\sidenote{The exact prefactor depends on the signaling scheme (voltage-mode,
current-mode, terminated transmission line) but the \emph{length dependence} is
the first principle that matters here.} So the energy to move a bit
electrically \emph{rises with distance}, and the resistive-capacitive delay and
the equalization needed to fight it rise with both distance and data rate.

\section{The optical alternative: energy at the ends}

An optical link spends its energy differently. The dominant costs are the
\emph{conversions at the two ends} (driving the laser or modulator, and the
receiver that turns light back into charge) plus the wall-plug inefficiency of
the laser. Over short reaches, waveguide and fiber loss are small, so the energy
per bit is, to first order, \emph{independent of distance}.

\keyidea{Electrical energy per bit grows with length; optical energy per bit is
set at the endpoints and is roughly flat with length. That single contrast is the
seed of everything else.}

\section{The break-even distance, and why optics moves toward the chip}

Put the two together and there is a cross-over length: below it, electrical wins;
above it, optical wins. The decisive observation is what happens as data rates
climb: the electrical cost at a given length rises (more charging events per
second, worse loss at higher frequency, more equalization), so the \emph{break-even
distance shrinks}.

That is the whole history of the field in one sentence. As rates went from
gigabits to hundreds of gigabits per lane, the cross-over marched inward: campus,
then rack, then board, then package, and now die-to-die. It is the first-principles
reason co-packaged optics and in-package optical I/O exist
(\Cref{sec:cpo-status,sec:cwwdm}), and the reason the scope of this book is the
\emph{shortest} links (\Cref{sec:reach}).

\section{The receiver, capacitance, and attojoule targets}

Miller's sharpest point concerns the receiver. To register a bit, the photocurrent
must develop a detectable voltage on the receiver's input node, so the detection
energy again looks like a $C V^{2}$ on that node's capacitance. Minimize the
photodetector and input capacitance (by integrating the detector tightly with the
first transistor, eliminating parasitic pads and wires) and you can detect a bit
with \emph{fewer photons and less energy}. This is the argument for close
electronic--photonic integration, and it is what makes sub-100~fJ/bit (and, in
principle, attojoule-class) devices conceivable.

\aside{This is why integration is not a packaging convenience but an energy
strategy: shortening the electrical path between detector and amplifier lowers
capacitance, which lowers both the energy and the optical power a link needs.}

\section{The floors: photons per bit and noise}

Two physical floors bound how far this can go. A real receiver needs enough
photons per bit for adequate signal-to-noise: today hundreds, with ideal detection
in the handful-of-photons range. That sets a floor on received optical power, and
hence on laser energy. Separately, the $kT/C$ noise on the receiver node (not the
Landauer $kT\ln 2$ limit of logic, which is far smaller) is the practical noise
floor for communication, and it too rewards small capacitance.

Miller is careful to separate the energy of \emph{logic} from the energy of
\emph{communication}; short-reach interconnect is overwhelmingly a communication-%
energy problem, dominated by the endpoints described above.

\section{Recent progress toward the limits}
\label{sec:miller-recent}

Miller's numbers are years old, but the framework has aged well: the last few
years have been a steady march of experiments toward the floors it predicts, and
they validate rather than overturn it. Three threads are worth knowing.

First, \textbf{low-capacitance 3D integration is delivering the receiver energy
Miller argued for}. A 2025 demonstration integrated an 80-channel transceiver in
three dimensions, stacking the photonics directly on the CMOS to minimize the
receiver-node capacitance, and reported roughly \SI{120}{fJ/bit}.\sidenote{\citehref{daudlin2025}{S.~Daudlin
\emph{et al.}, ``Three-dimensional photonic integration for ultra-low-energy,
high-bandwidth interchip data links,'' \emph{Nature Photonics} 19:502--509,
2025}.} Tellingly, it reaches that number not by pushing per-lane rate but by using
\emph{many slow channels} (about \SI{10}{Gb/s} each) so each receiver stays in
its most sensitive, lowest-energy regime, with WDM providing the aggregate
bandwidth. That is Miller's prescription almost verbatim.

Second, \textbf{the capacitance argument is directly measurable}. A co-designed
12\,nm-FinFET-on-silicon-photonics transceiver using direct-bond interconnect cut
input parasitic capacitance by about \SI{75}{\percent}, which bought
$\sim\!\SI{6}{dB}$ of receiver sensitivity and pushed link energy to a few hundred
fJ/bit.\sidenote{\citehref{chang2023}{P.-H.~Chang \emph{et al.}, ``A 3D integrated energy-efficient
transceiver \ldots{} co-designed 12\,nm FinFET and silicon photonic ICs,''
\emph{JLT} 41(21):6741--6755, 2023}.} The $C V^{2}$ story is not a metaphor; you can
watch the sensitivity move as the capacitance drops.

Third, \textbf{WDM receivers with near-zero-power wavelength control are
approaching sub-pJ/bit at terabit scale}: a single-chip 32-channel WDM PAM4
receiver reached \SI{1.024}{Tb/s} on one fiber at under \SI{0.38}{pJ/bit} with no
DSP or equalization.\sidenote{\citehref{pirmoradi2025}{A.~Pirmoradi \emph{et al.}, ``A single-chip
\SI{1.024}{Tb/s} silicon photonics PAM4 receiver,'' arXiv:2507.12452, 2025}.} And
on the device side, alternative emitters are being pushed into the same regime, a
2025 microLED link demonstrated \SI{200}{fJ/bit} transmitter energy at
$\text{BER} < 10^{-12}$ with no FEC,\sidenote{\citehref{avicena2025}{Avicena, ECOC 2025 demonstration of
the LightBundle microLED interconnect}.} exactly the LED branch of Miller's
device-by-device comparison.

\keyidea{The 2009/2017 framework has not been superseded; it has been confirmed in
hardware. The winning recipe in every recent result (low-capacitance
co-/3D-integration plus many WDM channels at modest per-lane rate) is precisely
what Miller's energy accounting recommends. What has changed is only the vertical
axis: from theoretical attojoule targets toward demonstrated hundreds-of-fJ/bit
links.}

\section{Why this framework anchors the book}

Everything that follows is an effort to approach these floors at the required data
rate and reliability. Laser wall-plug efficiency (\Cref{ch:lasers}) sets how much
optical power you can afford. Modulation and FEC (\Cref{ch:imdd,sec:kp4,sec:equalization})
trade SNR for reach and rate. WDM (\Cref{ch:wdm}) amortizes one laser source across
many wavelengths. Receiver noise and sensitivity
(\Cref{ch:models,sec:worked-budget,sec:link-budget}) decide whether that power is
enough. Co-packaging and energy-per-bit trends (\Cref{sec:power}) show the same
first principle playing out as copper reach collapses and optics move onto the
interposer.

\keyidea{Short-reach optics is, at bottom, an energy-per-bit optimization. Optical
energy is spent at the endpoints and is flat with distance; electrical energy grows
with distance and rate; so rising data rates push optics ever closer to the chip.
Miller's framework is the lens the rest of this book looks through.}
